% ============================================================================
% Requires in the preamble:
%   \usepackage{amsmath,amssymb,amsthm}
%   \newtheorem{assumption}{Assumption}
%   \newtheorem{definition}{Definition}
%   \newtheorem{remark}{Remark}
% ============================================================================

\section{Problem Formulation}
\label{preliminary}

\subsection{Preliminaries}

\paragraph{MDPs and the role of the dynamics.}
We model each domain as a Markov decision process (MDP)
$\mathcal{M} = (\mathcal{S}, \mathcal{A}, p, r, \gamma, \rho_0)$ with state space
$\mathcal{S}$, action space $\mathcal{A}$, transition kernel
$p(s' \mid s, a)$, reward function $r$, discount $\gamma \in [0,1)$ and initial
state distribution $\rho_0$. Among the components of this tuple, the transition
kernel $p$ --- the \emph{dynamics} --- is the only one that encodes physics: it
is where inertia, friction, contact stiffness and actuator response enter the
problem. The remaining components specify the task ($r$, $\gamma$, $\rho_0$) or
the agent's interface to it ($\mathcal{S}$, $\mathcal{A}$). Our entire setting
is built on this decomposition: we hold the task fixed and perturb only $p$.

\paragraph{Policies and values are solutions \emph{relative to} a kernel.}
A policy $\pi(a \mid s)$ interacting with $p$ induces the discounted return, the
value functions, and the normalized state--action occupancy
\begin{equation}
J_{p}(\pi) = \mathbb{E}_{\pi,p}\!\left[\sum_{t\ge0} \gamma^t r_t\right],
\quad
V^{\pi}_{p}(s) = \mathbb{E}_{\pi,p}\!\left[\sum_{t\ge0} \gamma^{t} r_t \,\Big|\, s_0 = s\right],
\quad
Q^{\pi}_{p}(s,a) = \mathbb{E}_{\pi,p}\!\left[\sum_{t\ge0} \gamma^{t} r_t \,\Big|\, s_0 = s, a_0 = a\right],
\end{equation}
\begin{equation}
d^{\pi}_{p}(s,a) = (1-\gamma)\sum_{t\ge0}\gamma^{t}\,
\Pr_{\pi,p}\!\left(s_t = s,\, a_t = a\right).
\label{eq:occupancy}
\end{equation}
We deliberately carry the kernel in the subscript, because the central
difficulty of this paper is that a policy is a \emph{kernel-independent object}
but a \emph{kernel-dependent solution}: the map $s \mapsto a$ can be executed
verbatim in any domain sharing $\mathcal{S}$ and $\mathcal{A}$, yet the property
that makes it worth executing --- optimality, or even competence --- was
established under one particular $p$ and does not travel with it.

The same is true of value functions, but one step more indirectly. $Q^{\pi}_{p}$
is defined only implicitly, as the fixed point of the Bellman operator
\begin{equation}
\left(\mathcal{T}_{p}\,Q\right)(s,a) \;=\; r(s,a) \;+\;
\gamma\, \mathbb{E}_{s' \sim p(\cdot \mid s,a)}\!\left[\,Q\!\left(s', \pi(s')\right)\right],
\label{eq:bellman}
\end{equation}
which is itself an expectation \emph{under $p$}. Value learning --- fitting
$Q$ by repeatedly applying a sampled version of \eqref{eq:bellman} to collected
transitions --- therefore never observes $p$ directly; it only ever sees $p$
through the successor states in its data. A critic trained to satisfy
$Q = \mathcal{T}_{p_s} Q$ is, when the kernel changes, simply no longer a fixed
point: its Bellman residual under the new kernel is the value-side signature of
a dynamics gap. Finally, note that the occupancy \eqref{eq:occupancy} depends on
\emph{both} arguments: the same $\pi$ visits different state distributions under
different kernels. This is why, in an adaptation setting, ``on-policy'' is not a
well-formed notion until one also says \emph{in which domain} --- a distinction
that will turn out to decide whether adaptation works at all.

\paragraph{Learned dynamics models.}
The true kernel is never available to the agent; a model-based method instead
fits a parametric approximation $\hat{p}_\theta \approx p$ from observed
transitions and uses it as a stand-in for $p$. Following recent latent
world-model architectures, we do not estimate a density over raw states but a
deterministic transition in a learned latent space: an encoder
$z = h_\phi(s)$ and a latent dynamics head $z' = d_\theta(z,a)$, trained by
multi-step latent consistency
\begin{equation}
\mathcal{L}_{\mathrm{dyn}}(\theta,\phi) \;=\;
\mathbb{E}_{(s_{0:H},\,a_{0:H-1}) \sim \mu}
\sum_{t=0}^{H-1} \lambda^{t}\,
\big\|\, d_\theta(\hat{z}_t, a_t) \;-\; \mathrm{sg}\!\left(h_\phi(s_{t+1})\right)
\big\|^{2},
\qquad \hat{z}_0 = h_\phi(s_0),\;\; \hat{z}_{t+1} = d_\theta(\hat{z}_t,a_t),
\label{eq:consistency}
\end{equation}
where $\mathrm{sg}$ denotes stop-gradient and $\lambda$ discounts longer
horizons. Two features of \eqref{eq:consistency} matter for what follows.
First, the model is trained by \emph{rolling its own predictions}, so its
accuracy is a property of $H$-step trajectories rather than of single
transitions. Second, and more consequentially, the objective is an expectation
under a sampling distribution $\mu$: a learned model is only constrained where
$\mu$ places mass. We will argue that in adaptation the choice of $\mu$ --- not
the capacity of $d_\theta$ or the size of the dataset --- is what determines
success.

\paragraph{Source and target domains.}
We instantiate two MDPs that share $(\mathcal{S}, \mathcal{A}, r, \gamma,
\rho_0)$ and differ only in their kernels: a \emph{source} domain with $p_s$,
in which interaction is cheap and unlimited, and a \emph{target} domain with
$p_t$, in which it is scarce. We abbreviate $J_s \!:=\! J_{p_s}$,
$J_t \!:=\! J_{p_t}$, and likewise for $V$, $Q$ and $d^\pi$. Unrestricted
reinforcement learning in the source yields an expert $\pi_s$ together with its
critic $Q^{\pi_s}_{p_s}$; deployment, however, happens under $p_t$. Since
$\pi_s$ was optimized against $\mathcal{T}_{p_s}$ and the contact-phase physics
it presumes, executing it under $p_t$ generally incurs a performance drop
\begin{equation}
\mathrm{gap}(\pi_s) \;=\; J_{s}(\pi_s) \;-\; J_{t}(\pi_s) \;>\; 0,
\label{eq:gapdef}
\end{equation}
which we refer to as the \emph{dynamics gap} and formalize in
Section~\ref{sec:setting}.

This leaves the agent with three transferable assets, each carrying the source
kernel's imprint in a different way, and each repairable at a different price:
the \textbf{policy} $\pi_s$, which is executable anywhere but optimal only under
$p_s$; the \textbf{value} $Q^{\pi_s}_{p_s}$, which is a fixed point of the wrong
Bellman operator but may still \emph{rank} states correctly if the task is
unchanged; and the \textbf{model} $\hat{p}_\theta$, which is the only asset that
represents the mismatched quantity explicitly and can therefore be corrected by
target data directly, without waiting for reward to propagate. The question this
paper asks is which of the three a limited target-interaction budget should be
spent on.

\paragraph{Sparse, goal-defined reward.}
We study contact-rich manipulation with a terminal success reward,
\begin{equation}
r(s,a) = \mathbf{1}\!\left[s \in \mathcal{G}\right],
\label{eq:sparse}
\end{equation}
where the goal set $\mathcal{G} \subset \mathcal{S}$ is a geometric predicate on
the object configuration (e.g. the needle tip inside the ring). The episode
terminates on entry into $\mathcal{G}$, so $J(\pi)$ is a discounted success
rate and $V^{\pi}(s) \in [0,1]$ is (up to discounting) the probability that
$\pi$ eventually solves the task from $s$. This is the regime in which
credit assignment is hardest: no dense signal shapes the approach, and a
value function only becomes informative through long, goal-reaching
trajectories.

\paragraph{Three ways to consume interaction data.}
The methods we compare differ in what they fit with target-domain data and
where the resulting object is used.
\emph{(i) Model-free} methods fit a critic directly on collected transitions,
$\min_{Q} \mathbb{E}_{\mathcal{D}} \big[ (Q(s,a) - r - \gamma
Q^{-}(s', \pi(s')))^2 \big]$, and improve $\pi$ against it; under
\eqref{eq:sparse} this requires enough interaction for reward information to
propagate backwards along the horizon.
\emph{(ii) Dyna}-style methods fit an approximate model $\hat{p}$ and use
\emph{synthetic} rollouts from $\hat{p}$ to supply extra transitions to the
same value/policy update, trading model bias for sample efficiency.
\emph{(iii) Planning} methods use $\hat{p}$ at decision time: given a learned
latent encoder $z = h_\phi(s)$, latent dynamics $z' = d_\theta(z,a)$, reward
head $\hat{r}$, terminal value head $\hat{V}$, and a proposal policy
$\hat{\pi}$, the executed action is the first element of
\begin{equation}
a_{0:H-1}^{\star} \;=\; \arg\max_{a_{0:H-1}}
\;\; \underbrace{\sum_{t=0}^{H-1} \gamma^{t}\, \hat{r}(z_t, a_t)}_{\text{within horizon}}
\;+\; \underbrace{\gamma^{H}\, \hat{V}(z_H)}_{\text{terminal bootstrap}},
\qquad z_{t+1} = d_\theta(z_t, a_t),
\label{eq:mppi}
\end{equation}
solved approximately by sampling-based optimization (MPPI/CEM) seeded from
$\hat{\pi}$. Crucially, \eqref{eq:mppi} consumes the value function only
through the \emph{ranking} it induces over the candidate terminal latents
$\{z_H^{(i)}\}_i$ that the model itself reaches, whereas it consumes the
dynamics $H$ times per candidate. Section~\ref{empirical} shows this
asymmetry is what decides which component must be adapted.

\subsection{Sim-to-real adaptation as a dynamics gap}
\label{sec:setting}

We now state precisely what the source and target domains share and what they
do not. Four assumptions define the setting; the first two are constructive
(they say how we build the gap), the last two are substantive (they say what we
believe about it, and both are tested in Section~\ref{empirical}).

\begin{assumption}[Shared task, shifted dynamics]
\label{as:gap}
The two domains differ \emph{only} in their transition kernels:
\begin{equation}
\mathcal{M}_s = (\mathcal{S}, \mathcal{A}, p_s, r, \gamma, \rho_0),
\qquad
\mathcal{M}_t = (\mathcal{S}, \mathcal{A}, p_t, r, \gamma, \rho_0),
\qquad p_s \neq p_t .
\end{equation}
In particular the goal set $\mathcal{G}$, and hence the reward
\eqref{eq:sparse}, is identical: \emph{what} must be achieved does not change,
only \emph{how} the physics responds.
\end{assumption}

\begin{assumption}[Observable gap: no latent registration shift]
\label{as:observable}
The two domains share the same state space and the same observation map, and
the shift is entirely in $p$; there is no unobserved pose, calibration or
context variable. Both domains are therefore MDPs over the same
$\mathcal{S}$, and the gap is $p_t(\cdot \mid s,a) \neq p_s(\cdot \mid s,a)$
even under full observability.
\end{assumption}

Assumption~\ref{as:observable} is a deliberate narrowing of the usual
sim-to-real setting, and it is what makes the empirical study in
Section~\ref{empirical} interpretable. A large part of the sim-to-real
literature manufactures the gap by perturbing object \emph{poses} while hiding
the perturbation from the policy; that construction is really a POMDP, and
methods can succeed there by inferring a latent variable rather than by
correcting physics. We instead perturb only inertial, frictional, actuator and
clearance properties (Section~\ref{empirical}), so that no amount of state
estimation can close the gap and every method is forced to confront
$p_t$ itself.

\begin{assumption}[Localized gap]
\label{as:local}
Let $\Delta(s,a) = D_{\mathrm{TV}}\!\left(p_s(\cdot \mid s,a),\,
p_t(\cdot \mid s,a)\right)$. The gap is concentrated on a small,
contact-dependent subset of the state--action space: $\Delta(s,a) \approx 0$
in free space and is large only during contact, so
$\mathbb{E}_{d^{\pi}}[\Delta]$ is small while $\sup_{s,a} \Delta(s,a)$ is not.
\end{assumption}

Assumption~\ref{as:local} is the empirically dominant regime for manipulation:
reaching and transport transfer almost perfectly, and the entire failure mass
concentrates in the alignment/insertion phase. It has two consequences that
organize the rest of the paper. First, a source-trained expert does not fail
\emph{gracefully} --- a small average model gap produces a large closed-loop
drop, because the affected transitions are exactly the ones on which task
success is decided. Second, and less obviously, \emph{a model can only be
repaired where it is exercised}: since $\Delta$ is supported on a thin region
that is reached only by a competent policy, any adaptation objective of the
form $\min_\theta \mathbb{E}_{(s,a) \sim \mu}\|d_\theta(h(s),a) - h(s')\|^2$
is only useful when the sampling distribution $\mu$ matches the occupancy
$d^{\pi_{\text{deploy}}}_{t}$ of the policy that will actually be
deployed. We refer to this requirement as \emph{on-policy consistency}, and it
is the design principle behind our method (Section~\ref{simopd}).

\begin{assumption}[Preserved reachability / value transferability]
\label{as:value}
The goal remains reachable in the target from the support of $\rho_0$, and the
shift does not create states that are recoverable in the source but terminally
unrecoverable in the target. Consequently the source value function remains a
valid \emph{ordering} of target states: for the candidate sets encountered in
\eqref{eq:mppi}, $\arg\max$ under $V^{\pi_s}_{s}$ agrees with
$\arg\max$ under $V^{\pi_t^\star}_{t}$ up to a residual smaller
than the spread that planning must resolve.
\end{assumption}

Assumption~\ref{as:value} is stated explicitly because it is the hinge of our
analysis rather than a technicality: it predicts that the value function need
\emph{not} be adapted, a prediction we test directly in
Section~\ref{empirical}. It is also falsifiable --- a shift that changes
reachability (e.g. one that lets the fixture topple irrecoverably) violates it,
and we return to this boundary in Section~\ref{lim}.

\paragraph{Why the dynamics error, and not the value error, is the binding
constraint.} Under Assumptions~\ref{as:gap}--\ref{as:value}, the two error
sources in \eqref{eq:mppi} enter with very different multiplicities. A
per-step latent dynamics error $\varepsilon_d$ compounds along the planning
horizon, displacing the predicted terminal latent by $O(\sum_{t<H} L^{t}
\varepsilon_d)$ for a locally Lipschitz $d_\theta$, and --- because the
optimizer \emph{selects} the candidates that the model believes are best ---
this displacement is adversarially, not randomly, distributed. A value error
$\varepsilon_V$, in contrast, enters once, discounted by $\gamma^{H}$, and only
through its \emph{variation across candidates}: adding any function that is
constant over $\{z_H^{(i)}\}$ leaves the $\arg\max$ in \eqref{eq:mppi}
unchanged. Planning is therefore first-order sensitive to $\varepsilon_d$ and
only sensitive to the \emph{candidate-wise contrast} of $\varepsilon_V$, which
Assumption~\ref{as:value} asserts is small. Sections~\ref{empirical} and
\ref{simopd} give the empirical counterpart of this argument, including the
regime where it breaks.

\subsection{Adaptation under an asymmetric budget}

The setting is defined as much by its budget asymmetry as by its dynamics gap.
Source interaction is essentially free (it is a simulator), so we allow an
unrestricted source phase; target interaction is the scarce resource.

\begin{definition}[Adaptation budget]
\label{def:budget}
An adaptation method is given
(i) unrestricted access to $\mathcal{M}_s$, from which it may obtain an expert
policy $\pi_s$, its value function $V^{\pi_s}_{s}$, and arbitrary
source rollouts;
(ii) a small set of $n$ target demonstrations
$\mathcal{D}_t = \{\tau^{(i)}\}_{i=1}^{n}$, $n \le 20$; and
(iii) at most $N$ target environment transitions of online interaction.
No target-domain reward shaping, no privileged access to $p_t$ or to its
parameters, and no target expert are available.
\end{definition}

\begin{definition}[Adaptation problem]
Given the budget of Definition~\ref{def:budget}, output a deployable policy
$\pi$ maximizing $J_{t}(\pi)$. Methods are compared by
(a) the target success rate at a fixed budget $N$, and (b) the budget $N$
required to reach a fixed success level, i.e. their position on the
success-versus-interaction curve rather than at a single point.
\end{definition}

Both axes matter and they order methods differently, which is the subject of
Section~\ref{empirical}: the methods that are strongest asymptotically are not
the ones that move first. We stress that the regime of interest is the
\emph{pre-asymptotic} one. A sim-to-real practitioner operates at a budget
fixed by how much target interaction is affordable, which in our setting is
strictly less than what unassisted value learning needs; asymptotic performance
under an unbounded budget is therefore not the objective, and we invoke it only
as a diagnostic reference (Section~\ref{sec:mf-vs-mb}) to establish what a
given target domain \emph{admits}, so that a method's plateau can be attributed
either to the budget or to its own approximations.

\paragraph{Measuring the gap.}
We report the magnitude of a dynamics gap by the \emph{zero-shot drop} of the
frozen source expert, i.e. $\mathrm{gap}(\pi_s)$ of \eqref{eq:gapdef},
evaluated under the same success predicate in both domains. This is the
operationally meaningful measure --- it is exactly the performance the
adaptation phase must recover --- and, unlike a divergence between transition
kernels, it is comparable across tasks with different state
parameterizations.

\section{Empirical Studies of Adaptation to Dynamics Gap}
\label{empirical}

\subsection{Experimental setting}
\label{sec:setup}

\paragraph{Sim-to-sim as a controlled proxy.}
We instantiate Assumptions~\ref{as:gap}--\ref{as:observable} in simulation, by
holding the task, the asset and the success predicate fixed and perturbing only
the physical parameters of the MuJoCo model. This buys three things that a
physical sim-to-real study cannot provide, and all three are used in this
paper: (i) the gap is a \emph{known, dialable} quantity, so we can trace
dose--response curves rather than report a single transfer number;
(ii) an \emph{oracle dynamics} ablation is available --- we can run the exact
same planner with the true target simulator in place of the learned model,
which isolates model error from every other source of failure; and
(iii) evaluation is cheap enough to run hundreds of episodes per condition,
which is a prerequisite for reading effects of the size we study.

\paragraph{Source and target domains.}
The source domain is the nominal simulator. The target domain applies a fixed
perturbation to contact-phase physics: object mass and inertia, sliding
friction of the manipulated object and of the fixture, gripper actuator
stiffness, and the insertion clearance (the geometric aperture the held object
must pass). Two properties of the construction are deliberate. First, the
success predicate is defined by a fixed tolerance on the object pose and is
\emph{not} re-read from the perturbed model, so shrinking the physical
clearance adds a real obstacle without moving the goal --- the reward is
invariant as required by Assumption~\ref{as:gap}. Second, no positional offset
is applied to any object, so nothing about the shift is hidden from the state:
the gap is a pure dynamics gap in the fully observed MDP
(Assumption~\ref{as:observable}). We use three contact-rich manipulation tasks
of increasing tolerance width [TASKS], and calibrate each target so that the
frozen source expert drops from near-perfect source performance to a low but
nonzero target success rate, which keeps the setting an adaptation problem
rather than an exploration problem.

\paragraph{Source-domain training.}
In each task we train a source expert to near-saturation with a
demonstration-regularized off-policy RL method, and take from it both the
policy $\pi_s$ and its critic, from which the transferable value
$V^{\pi_s}_{s}(s) = \min_j Q_j(s, \pi_s(s))$ is read. All
adaptation methods start from this same source asset, so that comparisons
isolate what is done with target data.

\paragraph{Target budget control.}
Every method receives the same budget (Definition~\ref{def:budget}): the same
$n$ target demonstrations, collected once by a scripted/expert collector and
shared verbatim across arms, and the same cap on target transitions. Curves
are plotted against target transitions consumed, and model-free, Dyna and
planning arms are compared at matched budgets.

\paragraph{Evaluation protocol.}
Because the effects we study are of the order of a few success-rate points, we
fix the following protocol. All reported numbers use $\ge 300$ evaluation
episodes with seeds held fixed across arms; headline numbers use independent
re-evaluation of a saved checkpoint rather than the in-training evaluation,
which is systematically optimistic; and comparisons across checkpoints are only
made when the full environment configuration (shift parameters, initial-state
distribution, success predicate, planner hyper-parameters) is verified
identical. We report the noise floor of each comparison alongside the effect.

% ============================================================================
% NOTE (preamble additions needed for the floats below):
%   \usepackage{graphicx}
%   \usepackage{booktabs}
%   \usepackage{caption}      % for \captionof inside minipage
%   \usepackage{subcaption}   % for \subcaptionbox / subfigure
%   \usepackage{multirow}
%   \usepackage{wrapfig}      % for the wrapfigure in Sec.~\ref{sec:what-to-adapt}
%   \usepackage{xcolor}
%   \newcommand{\ph}[1]{\textcolor{gray}{#1}}   % placeholder numbers
%   (\shortstack and \underline are kernel; no package needed)
% ============================================================================

% ---------------------------------------------------------------------------
% TODO(global): task set. Threading_D0 / Coffee_D0 / Square_D0 are running;
% ThreePieceAssembly (TPA) is being tuned. Every float below is written for
% 4 tasks; if TPA does not land, delete the last column/panel and change
% "four tasks" -> "three tasks" in the text. Search for [TPA] to find them.
% TODO(global): every cell must carry the eval-episode count and be produced by
% independent re-evaluation (not model0.pt / in-training eval), per Sec.~\ref{sec:setup}.
% ---------------------------------------------------------------------------

\subsection{Model-Free vs.\ Model-Based Adaptation: why do we need a world model?}
\label{sec:mf-vs-mb}

[TODO: Clarify the trade-offs between model-based and model-free methods.
Model-based methods leverage a world model to generate additional data, yielding
higher sample efficiency; they are, however, hindered by the difficulty of
fitting an accurate model, and long-horizon rollouts in the latent space suffer
from compounding error. Model-free methods are simpler and more direct;
nevertheless, even off-policy algorithms need a large amount of interaction
before the value function converges, which under a sparse terminal reward
(Eq.~\ref{eq:sparse}) is precisely what a scarce target budget cannot buy.
IMPORTANT FRAMING: the question of this subsection is \emph{not} which method
wins with unlimited interaction. The budget is fixed by Def.~\ref{def:budget},
and the operating point lies entirely inside it. Model-free RL is therefore not
a competitor we lose to asymptotically; within the budget it is simply not yet
functional, and we run it out of budget only as a \emph{diagnostic reference}
that tells us what the target domain admits.]

% ---------------------------------------------------------------------------
% FLOAT 1 (Sec.~\ref{sec:mf-vs-mb}): table (left) + figure (right), one row.
%   LEFT  = budget-to-threshold + matched-budget score  -> "inside the budget,
%           only the model-based arm is functional at all"
%   RIGHT = DIAGNOSTIC, deliberately run past the budget -> "the Dyna plateau is
%           a model-error ceiling, not a data ceiling": something else reaches
%           near-saturation on the same target, so the plateau is the model's
%           fault and not the task's.
% The two halves are the two axes of Def.~\ref{def:budget}: (b) then (a).
% DO NOT frame the right half as "model-free eventually wins" -- an out-of-budget
% score is not a result in this setting, and reading it as a competitor would
% concede a comparison we never make. Its only job is attribution.
% ---------------------------------------------------------------------------
\begin{figure*}[t]
\centering
\begin{minipage}[t]{0.52\textwidth}
  \centering
  \captionof{table}{%
    \textbf{Sample efficiency of model-based adaptation.} Target transitions
    required to reach $30\%$ and $50\%$ success, and success rate at a matched
    budget of $100$k target transitions. Arrows indicate whether lower or
    higher is better; \texttt{---} means the threshold was never reached within
    the budget. All entries are means over [TODO: $k$] seeds with
    $\ge 300$ evaluation episodes; $\pm$ is the seed spread.
    % TODO: also report the zero-shot score of the frozen source expert
    % (Eq.~\ref{eq:gapdef}) as a "0 transitions" reference row -- it is the
    % starting point every arm shares, and it makes the speedup readable.
  }
  \label{tab:mf-vs-mb-budget}
  \resizebox{\linewidth}{!}{%
  \begin{tabular}{llccc}
  \toprule
   &  & \multicolumn{2}{c}{Steps to threshold $\downarrow$} & SR @100k $\uparrow$ \\
  \cmidrule(lr){3-4}\cmidrule(lr){5-5}
  Task & Method & $30\%$ & $50\%$ & \\
  \midrule
  \multirow{2}{*}{Threading}
    & IBRL (model-free) & \ph{--k} & \ph{--k} & \ph{0.--} \\
    & \textbf{Dyna-IBRL} & \ph{--k} & \ph{--k} & \ph{0.--} \\
  \midrule
  \multirow{2}{*}{Coffee}
    & IBRL (model-free) & \ph{--k} & \ph{--k} & \ph{0.--} \\
    & \textbf{Dyna-IBRL} & \ph{--k} & \ph{--k} & \ph{0.--} \\
  \midrule
  \multirow{2}{*}{Square}
    & IBRL (model-free) & \ph{--k} & \ph{--k} & \ph{0.--} \\
    & \textbf{Dyna-IBRL} & \ph{--k} & \ph{--k} & \ph{0.--} \\
  \midrule
  \multirow{2}{*}{TPA~[TPA]}
    & IBRL (model-free) & \ph{--k} & \ph{--k} & \ph{0.--} \\
    & \textbf{Dyna-IBRL} & \ph{--k} & \ph{--k} & \ph{0.--} \\
  \midrule
  \multicolumn{2}{l}{\emph{Speedup to $50\%$ (geo.\ mean)}} & & \ph{$--\times$} & \\
  \bottomrule
  \end{tabular}}
\end{minipage}\hfill
\begin{minipage}[t]{0.45\textwidth}
  \centering
  \includegraphics[width=\linewidth]{figs/placeholder_mf_vs_mb_asymptote.pdf}
  \caption{%
    \textbf{The model-based plateau is a model-error ceiling, not a budget
    ceiling.} \emph{Diagnostic, deliberately run outside the adaptation
    budget.} Target success rate vs.\ target transitions on [TODO: 2
    representative tasks, one panel each], continued to [TODO: $\sim$1M]
    transitions --- roughly [TODO: $--\times$] the budget of
    Def.~\ref{def:budget}. The shaded region on the left is the budget we
    actually operate in; inside it the model-free arm has barely left the
    zero-shot level, which is why it is not a viable adaptation method here.
    Outside it, model-free IBRL nonetheless climbs to $\ph{0.--}$, while
    Dyna-IBRL stops improving at $\ph{0.--}$ despite being given the same
    extended interaction. The comparison is used for attribution only: since the
    same target domain \emph{does} admit near-saturating performance, the
    model-based plateau cannot be blamed on the task or on the amount of data,
    and must be a property of the learned model --- which is what
    Sections~\ref{sec:dyna-vs-plan} and \ref{sec:what-to-adapt} then localize.
    Shaded band = seed spread; horizontal dashed line = frozen source expert
    (zero-shot).
    % TODO(fig): SHADE the budget region and label it "adaptation budget".
    % The figure must read left-to-right as "in budget: nothing else works" ->
    % "out of budget: the plateau is the model's fault". Do not annotate the
    % crossover as an operating point: we always operate left of it.
    % TODO(fig): x-axis log scale, otherwise the in-budget region collapses to
    % a sliver and the early model-based advantage is invisible.
    % TODO(fig): give BOTH arms the same extended budget on the right. If the
    % Dyna arm is cut off at 100k the plateau claim is unsupported -- a curve
    % that was never run cannot be shown to have stopped rising.
    % TODO(fig): if the plateau is not flat, say so honestly -- do NOT crop the
    % x-range to manufacture one. A slowly-rising Dyna curve still supports the
    % efficiency claim on the left; it just weakens this attribution, in which
    % case fall back to the horizon sweep below as the primary evidence.
  }
  \label{fig:mf-vs-mb-asymptote}
\end{minipage}
\end{figure*}

[TODO: text for Table~\ref{tab:mf-vs-mb-budget}. Claim 1 --- \emph{under the
budget, the world model is not an optimization, it is what makes adaptation
possible at all}. Dyna-IBRL reaches $50\%$ in $\ph{--\times}$ fewer target
transitions than model-free IBRL, and at the matched $100$k budget it leads by
$\ph{+0.--}$ on average; on [TODO: which] tasks the model-free arm does not
reach the threshold within the budget at all. Attribute this to amortization:
each real transition is reused across many value updates through synthetic
rollouts, which is exactly the cost a sparse terminal reward
(Eq.~\ref{eq:sparse}) imposes on unassisted value learning. Phrase the takeaway
in budget terms, not in ranking terms --- ``$N_{50}$ falls inside the affordable
budget only for the model-based arm''.]

[TODO: text for Fig.~\ref{fig:mf-vs-mb-asymptote}. Claim 2 --- \emph{attribution
of the plateau}, stated explicitly as a diagnostic and not as a comparison.
Given $\ph{--\times}$ the budget, model-free IBRL reaches $\ph{0.--}$ while
Dyna-IBRL stops at $\ph{0.--}$ with the same extra interaction. Say plainly what
this does and does not license: it does \emph{not} say model-free RL is the
better adaptation method (at the budget of Def.~\ref{def:budget} it is not a
method at all), it says the model-based plateau is not caused by the budget.
Since the task, the target kernel and the success predicate are identical across
the two arms, the residual must come from the approximations the model-based arm
introduces --- and the rest of Section~\ref{empirical} is a search for which one.
% TODO(writing): pre-empt the obvious referee question here in one sentence ---
% "why not simply spend more target interaction?" -- by pointing back to the
% budget asymmetry of Sec.~\ref{def:budget}: in the sim2real problem this proxy
% stands in for, the right-hand side of this plot is not purchasable.]

% TODO(analysis, optional but strong): decompose the Dyna plateau by rollout
% horizon -- sweep the synthetic rollout length h and show the plateau height is
% non-monotone in h (short h = no value propagation under sparse reward, long h =
% compounding latent error). This is the bridge into Sec.~\ref{sec:dyna-vs-plan}
% and it is cheap: it reuses the same world model.

\subsection{Dyna vs.\ Planning: how should we use the world model for adaptation?}
\label{sec:dyna-vs-plan}

[TODO: Contrast Dyna with planning: planning significantly outperforms Dyna and
is substantially less vulnerable to dynamics error. Data generated in the latent
space by Dyna cannot be rolled out over long horizons without policy/model
collapse; yet under a sparse terminal reward, long trajectories reaching all the
way to the goal are exactly what value propagation needs. Dyna is thus squeezed
from both sides, and its three-way objective (dynamics, value and policy
optimized jointly) is notoriously hard to converge. Planning only has to keep
the dynamics honest during adaptation (e.g.\ SimDist adapts the dynamics head
alone), which is more direct, more tractable and more robust.]

% ---------------------------------------------------------------------------
% FLOAT 2 (Sec.~\ref{sec:dyna-vs-plan}): one wide figure, 3--4 panels,
% one panel per task. Matched-budget adaptation curves.
% ---------------------------------------------------------------------------
\begin{figure*}[t]
\centering
\begin{subfigure}[t]{0.245\textwidth}
  \includegraphics[width=\linewidth]{figs/placeholder_dvp_threading.pdf}
  \caption{Threading}\label{fig:dvp-threading}
\end{subfigure}\hfill
\begin{subfigure}[t]{0.245\textwidth}
  \includegraphics[width=\linewidth]{figs/placeholder_dvp_coffee.pdf}
  \caption{Coffee}\label{fig:dvp-coffee}
\end{subfigure}\hfill
\begin{subfigure}[t]{0.245\textwidth}
  \includegraphics[width=\linewidth]{figs/placeholder_dvp_square.pdf}
  \caption{Square}\label{fig:dvp-square}
\end{subfigure}\hfill
\begin{subfigure}[t]{0.245\textwidth}
  \includegraphics[width=\linewidth]{figs/placeholder_dvp_tpa.pdf}
  \caption{TPA~[TPA]}\label{fig:dvp-tpa}
\end{subfigure}
\caption{%
  \textbf{At a matched target budget, planning adapts far better than Dyna.}
  Target success rate vs.\ target transitions consumed (cap $100$k) for
  Dyna-IBRL and planning-based adaptation (SimDist), both built on the
  \emph{same} source-distilled world model and the same $n$ target
  demonstrations, so the comparison isolates \emph{how} the model is used ---
  as a data generator or as a decision-time simulator. Model-free IBRL at the
  same budget and the frozen source expert are shown as references; the shaded
  band is the seed spread and the grey band at the right edge marks the
  evaluation noise floor.
  % TODO(fig): the anchor number to reproduce here is the same-world-model
  % contrast in docs/on-policy-adaptation.md (planning 0.813 vs Dyna 0.157).
  % It MUST be the identical world-model checkpoint in both arms, otherwise the
  % comparison silently becomes a model-quality comparison.
  % TODO(fig): keep the in-budget model-free curve on the same axes -- it is the
  % "what a budget buys without a world model" reference, consistent with the
  % left half of Fig.~\ref{fig:mf-vs-mb-budget}. Do NOT extend it past the
  % budget cap here; the out-of-budget continuation lives only in the diagnostic
  % Fig.~\ref{fig:mf-vs-mb-asymptote}.
  % TODO(fig): compare only steady-state segments; the slower arm shows
  % transient peaks (methodology, Sec.~\ref{sec:setup}).
}
\label{fig:dyna-vs-planning}
\end{figure*}

[TODO: text for Fig.~\ref{fig:dyna-vs-planning}. Report the matched-budget
margin (planning $\ph{0.--}$ vs.\ Dyna $\ph{0.--}$, averaged over tasks) and
make the mechanism claim: with the same model, planning re-queries the model
$H$ steps ahead from the \emph{current} state and never has to commit model
error into learned parameters, while Dyna writes model error into the critic,
where it is compounded by bootstrapping and cannot be undone by later real data.]

[TODO: close the subsection with the honest limitation, so that
Sec.~\ref{sec:what-to-adapt} has a reason to exist: planning is the best use of
the budget, yet it saturates at $\ph{0.--}$ and additional target transitions
inside a realistic budget do not move it. Measure that residual against what is
\emph{achievable} on this target --- the oracle-dynamics planner and the target
expert (Sec.~\ref{sec:what-to-adapt}) --- rather than against the out-of-budget
model-free run, which is a diagnostic and not a target to catch. The remaining
question is which component of the planner the residual lives in.]

\subsection{What to adapt: value, dynamics, or policy?}
\label{sec:what-to-adapt}

[TODO: Frame the intuition and then refute it. One might suspect the ceiling of
planning is set by the value function, since a source--target value gap
demonstrably exists. Our experiments show the opposite: every attempt to adapt
the value degrades performance. The reason is that a value function is harder to
fit than a dynamics model --- and adapting it forfeits the very advantage of
planning, because like Dyna and model-free RL it needs reward to propagate, and
therefore needs a budget we do not have.]

% ---------------------------------------------------------------------------
% FLOAT 3 (Sec.~\ref{sec:what-to-adapt}, part A): "adapting the value breaks it"
% Grouped bars, one group per task, five arms. Includes the *marginal* value
% tweak so the refutation is not only about aggressive retraining.
% ---------------------------------------------------------------------------
\begin{figure*}[t]
\centering
\includegraphics[width=\textwidth]{figs/placeholder_value_adapt_bars.pdf}
\caption{%
  \textbf{Adapting the value function degrades planning, in every ordering and
  at every strength.} Final target success rate for five adaptation schedules
  under an identical target budget:
  (i) \emph{dynamics only} (SimDist, ours);
  (ii) \emph{joint} --- value and dynamics adapted simultaneously;
  (iii) \emph{value $\rightarrow$ dynamics};
  (iv) \emph{dynamics $\rightarrow$ value};
  (v) \emph{marginal value} --- the value head nudged with a small
  learning rate / few steps / small trust region [TODO: name the exact knob].
  Dynamics-only is best on all [TODO: four] tasks; no ordering rescues value
  adaptation, and the marginal variant merely interpolates towards the failure
  rather than avoiding it. Error bars are the seed spread and the grey band is
  the evaluation noise floor, so bars inside it are not distinguishable.
  % TODO(fig): the mechanism number belongs in the caption or the text -- the
  % candidate-discrimination signal available to the planner is
  % std_i(gamma^H V) = 0.0068, i.e. 1.4% of the value range, so *any* offline
  % refit residual is of the same order and rewrites the elite set.
  % Cross-ref the first-order argument in Sec.~\ref{sec:setting}.
  % TODO(fig): include the dose--response version as an inset or an appendix
  % panel: sweep the value-adaptation strength and show the score decreases
  % monotonically in it (Spearman = -1.000). Monotonicity is what turns
  % "our value adaptation was badly tuned" into "the direction is wrong".
  % TODO(fig): add ONE caveat arm if space allows -- the same sweep on a
  % *more accurate* dynamics model, where the degradation shrinks by ~9x.
  % That is the boundary of the claim: value can be touched only in proportion
  % to how good the dynamics already is.
}
\label{fig:value-adapt-bars}
\end{figure*}

% ---------------------------------------------------------------------------
% FLOAT 4 (Sec.~\ref{sec:what-to-adapt}, part B): the factorial that localizes
% the bottleneck. dynamics {oracle, learned} x value {source-distilled, target}
% x proposal {source policy, target BC (not fine-tuned)}.
% Two claims come out of it: (1) dynamics accuracy dominates, (2) when the
% dynamics is imperfect, the three components must be trained on ONE distribution.
% ---------------------------------------------------------------------------
\begin{table*}[t]
\centering
\caption{%
  \textbf{Localizing the bottleneck: dynamics accuracy dominates, and value is
  not the constraint.} Target success rate of the \emph{same} MPPI controller
  (Eq.~\ref{eq:mppi}) as we independently swap its three inputs: the dynamics
  (learned latent model vs.\ the true target simulator as an oracle), the
  terminal value (source-distilled $\hat V$ vs.\ a value fit on target data),
  and the proposal policy (frozen source expert $\pi_s$ vs.\ behaviour cloning
  on the $n$ target demonstrations, \emph{without} any further fine-tuning).
  [TODO: state the effect sizes as main effects: dynamics
  $\ph{+0.--}$ ($z=\ph{--}$) $\gg$ proposal $\ph{+0.--}$ $>$ value
  $\ph{+0.--}$ (n.s.).] All cells: [TODO: $\ge 300$] episodes, [TODO: $k$] seeds.
  % TODO(tab): report effects BOTH additively and as "fraction of the remaining
  % gap closed" (or log-odds). Near the ceiling an additive null is an
  % arithmetic artefact, not evidence that the mechanism disappeared.
  % TODO(tab): every row must be verified to share the identical env config
  % (sigma, init distribution, success predicate, planner hyper-parameters);
  % a mismatched pair here previously produced a wholly spurious conclusion.
  % TODO(tab): planner hyper-parameters do not transfer across tasks nor across
  % oracle/learned dynamics -- each row needs its own re-tuned init_std /
  % base_policy_std, otherwise this table measures tuning, not mechanism.
}
\label{tab:factorial-bottleneck}
\resizebox{\textwidth}{!}{%
\begin{tabular}{lllcccc}
\toprule
Dynamics & Value & Proposal & Threading & Coffee & Square & TPA~[TPA] \\
\midrule
Oracle  & source-distilled $\hat V$ & source $\pi_s$          & \ph{0.--} & \ph{0.--} & \ph{0.--} & \ph{0.--} \\
Oracle  & source-distilled $\hat V$ & target BC               & \ph{0.--} & \ph{0.--} & \ph{0.--} & \ph{0.--} \\
Oracle  & target-fit $V_t$          & source $\pi_s$          & \ph{0.--} & \ph{0.--} & \ph{0.--} & \ph{0.--} \\
Oracle  & target-fit $V_t$          & target BC               & \ph{0.--} & \ph{0.--} & \ph{0.--} & \ph{0.--} \\
\midrule
Learned & source-distilled $\hat V$ & source $\pi_s$          & \ph{0.--} & \ph{0.--} & \ph{0.--} & \ph{0.--} \\
Learned & source-distilled $\hat V$ & target BC               & \ph{0.--} & \ph{0.--} & \ph{0.--} & \ph{0.--} \\
Learned & target-fit $V_t$          & source $\pi_s$          & \ph{0.--} & \ph{0.--} & \ph{0.--} & \ph{0.--} \\
Learned & target-fit $V_t$          & target BC               & \ph{0.--} & \ph{0.--} & \ph{0.--} & \ph{0.--} \\
\midrule
\multicolumn{3}{l}{\emph{Main effect: dynamics (oracle $-$ learned)}}
  & \ph{+0.--} & \ph{+0.--} & \ph{+0.--} & \ph{+0.--} \\
\multicolumn{3}{l}{\emph{Main effect: value (target $-$ source)}}
  & \ph{+0.--} & \ph{+0.--} & \ph{+0.--} & \ph{+0.--} \\
\multicolumn{3}{l}{\emph{Main effect: proposal (target BC $-$ source $\pi_s$)}}
  & \ph{+0.--} & \ph{+0.--} & \ph{+0.--} & \ph{+0.--} \\
\bottomrule
\end{tabular}}
\end{table*}

[TODO: text for Table~\ref{tab:factorial-bottleneck}, claim A --- \emph{the
value is not the bottleneck}. The top-left cell (oracle dynamics $+$
source-distilled value $+$ source policy proposal) already reaches
$\ph{0.--}$, i.e.\ near-saturation, \emph{without a single target-domain value
update}. This is the direct empirical test of Assumption~\ref{as:value}: the
source value ranks the candidate terminal latents well enough, and what the
planner was missing was never the value.]

[TODO: text, claim B --- \emph{dynamics accuracy is the binding constraint}.
Holding value and proposal fixed, replacing the learned model with the oracle is
worth $\ph{+0.--}$, an order of magnitude more than either other swap, and it
alone accounts for the residual gap to model-free RL identified in
Sec.~\ref{sec:dyna-vs-plan}. Tie back to the multiplicity argument in
Sec.~\ref{sec:setting}: the dynamics is consumed $H$ times per candidate and
its error is selected \emph{adversarially} by the optimizer, while the value is
consumed once and only through candidate-wise contrast.]

[TODO: text, claim C --- \emph{when the dynamics is imperfect, the three
components must live on one distribution}. Read the interaction, not the main
effects: the value swap is harmless (or helps) under the oracle but
$\ph{--}$ under the learned model, and the effect of the proposal swap is
conditional on which value is used. Formulate this as the on-policy-consistency
statement of Sec.~\ref{sec:setting}: the model is only trustworthy on the
occupancy it was fit on, so a proposal or a value that moves the search off that
occupancy is punished in proportion to the model's error. This is the design
principle carried into Sec.~\ref{simopd}. Also flag the falsified alternative:
under an oracle the ``value must be self-consistent with the continuation
policy'' story does not hold, so the constraint is a property of model error,
not of value coherence per se.]

% ---------------------------------------------------------------------------
% FLOAT 5 (Sec.~\ref{sec:what-to-adapt}, part C): the policy/proposal axis.
% Small single-column curve: adaptation with vs. without a target-anchored
% proposal warm-up. Motivates the source-anchored fine-tuning in Sec.~\ref{simopd}.
% ---------------------------------------------------------------------------
% NOTE: wrapfigure must be issued at the START of the paragraph it wraps into,
% must not be adjacent to another float / list / section break, and does not
% survive a page break -- if it misbehaves at camera-ready, fall back to
% \begin{figure}[t] ... \end{figure}. Width below is the wrap column; the
% {r} places it on the right and {12} is the number of wrapped lines (retune
% once the surrounding paragraph is written).
\begin{wrapfigure}{r}{0.46\textwidth}
\vspace{-\baselineskip}
\centering
\includegraphics[width=\linewidth]{figs/placeholder_policy_warmup.pdf}
\caption{%
  \textbf{The proposal policy must be warmed up on target data --- but anchored
  to the source.} Target success rate over the course of adaptation, with the
  planner's proposal (a) frozen at the source expert, (b) fine-tuned on the $n$
  target demonstrations \emph{without} regularization, and (c) fine-tuned with a
  source anchor [TODO: name the coefficient]. Warm-up decides the level the
  whole adaptation starts from and converges to ($\ph{+0.--}$ for (c) over (a)),
  while unregularized fine-tuning (b) destroys the proposal --- $\ph{--}$
  demonstration transitions are re-visited for $\ph{--}$ effective epochs, so
  the proposal leaves the region where the model and value are valid.
  % TODO(fig): show ALL of (a),(b),(c) -- (b) is what makes the anchor a
  % mechanism claim rather than a hyper-parameter. Note that the size of (b)'s
  % damage is target-dependent (observed to vary by >20x across targets), so
  % report per-task and do not average it away.
  % TODO(fig): this panel is the empirical bridge to Sec.~\ref{simopd}: the
  % anchor is what keeps the proposal on the occupancy the dynamics was fit on,
  % i.e. it is the same on-policy-consistency principle as claim C above.
  % TODO(fig): as a wrapfigure this panel is ~0.46\textwidth -- redraw it for
  % that aspect ratio (taller than wide, legend inside the axes or dropped in
  % favour of direct line labels, 3 curves max). Do not shrink the full-width
  % version; the tick labels will be unreadable.
  % TODO(fig): if the wrap still crowds the text, merge it as a 5th group of
  % bars into Fig.~\ref{fig:value-adapt-bars}; keep it separate only if the
  % *curve shape* (starting level vs. rate of improvement) is part of the
  % argument -- which it is, since the claim is about where adaptation starts
  % AND where it converges.
}
\label{fig:policy-warmup}
\vspace{-\baselineskip}
\end{wrapfigure}

[TODO: text for Fig.~\ref{fig:policy-warmup} and the summary of the subsection.
State the three conclusions in the order they were established: (1) do not adapt
the value --- the planner consumes it only as a ranking and the available
candidate-wise signal is smaller than any offline refit residual; (2) improving
the dynamics accuracy is the one intervention that moves the ceiling; (3) while
the dynamics remains imperfect, the dynamics, the value and the proposal must be
trained and used on a single, shared occupancy, which is what the
source-anchored proposal warm-up buys. These three points define the method in
Sec.~\ref{simopd}.]

\section{SimOPD: Sim-to-real On-policy Adaptation with Ensemble Dynamics}
\label{simopd}

[TODO: Next, we present our key improvements: 1. Dynamics ensemble with
optimistic exploration and pessimistic planning. 2. Addressing the trade-off in
target domain demonstration fine-tuning for the policy proposal: We observe that
a prerequisite for MPPI to function effectively is that the policy proposal must
be trained on the same distribution as the value function. However, the policy
also requires target domain information to warm up and achieve superior
performance. To balance this, we introduce a KL divergence regularizer alongside
the behavioral cloning (BC) objective to constrain the fine-tuned policy from
drifting too far from the original policy, which significantly enhances
performance.]

% ---------------------------------------------------------------------------
% FLOAT 6 (Sec.~\ref{simopd}): MAIN RESULT.
% One wide figure, one panel per task (4 manipulation + 1 IsaacLab locomotion).
% Each panel: target success rate (or normalized return, for locomotion) vs.
% target transitions consumed, x-axis capped at the budget of Def.~\ref{def:budget}.
% Curves are grouped by FAMILY, and the families must be visually separable at a
% glance -- that is the whole point of using curves instead of a table:
%   model-free  -> cool colours, dotted
%   Dyna        -> warm colours, dashed
%   planning    -> solid, SimDist in a muted solid, SimOPD in the accent colour
%                  and the only curve with a heavier line weight
%   references  -> horizontal lines, no markers: frozen pi_s (grey dash-dot,
%                  the y-intercept every arm starts from) and the
%                  oracle-dynamics planner (black dotted, the achievable ceiling)
%
% WHY curves and not a final-score table: the claim of this paper is about the
% *budget axis*, so the evidence should live on that axis. A curve shows three
% things a single number cannot -- where an arm starts, how fast it moves, and
% whether it has stopped moving by the cap. In particular the model-free arms'
% flatness inside the budget is the argument of Sec.~\ref{sec:mf-vs-mb}, and it
% is invisible in an endpoint table.
%
% TODO(fig): decide the arm list; keep it to <=7 curves per panel or the panels
% become unreadable at column width. Proposed:
%   [mf]  IBRL fine-tuned from pi_s          (strongest model-free arm we have)
%   [mf]  residual RL on frozen pi_s         ("don't touch the expert" family)
%   [mf]  RLPD / demo-mixed off-policy RL    (standard demo-regularized recipe)
%   [dyna] Dyna-IBRL
%   [plan] TD-MPC2-style online world model
%   [plan] SimDist  (single dynamics head, no proposal anchor)  <- runner-up
%   [ours] SimOPD
% Two zero-online-interaction baselines (target BC on the n demos; a
% domain-randomized source expert) have no curve -- draw them as extra
% horizontal reference lines, or as points at x=0, and say which in the caption.
% Referees WILL ask for the DR baseline, so do not drop it just because it is
% not a curve.
%
% TODO(fig): every arm gets the SAME n demos (verbatim, same files) and the SAME
% transition cap. Arms that cannot consume the online half of the budget must be
% marked as such, not silently credited with it.
% TODO(fig): x-axis STOPS at the budget cap. The out-of-budget model-free
% continuation exists only in the diagnostic Fig.~\ref{fig:mf-vs-mb-asymptote};
% mixing budgets inside the main result is the easiest way to lose the framing.
% TODO(fig): shade +-1 seed spread; state n_seeds and eval episodes in the caption.
% Do not plot in-training eval -- it is systematically optimistic; curves must be
% built from periodic independent re-evaluation of saved checkpoints.
% ---------------------------------------------------------------------------
\begin{figure*}[t]
\centering
\begin{subfigure}[t]{0.195\textwidth}
  \includegraphics[width=\linewidth]{figs/placeholder_main_threading.pdf}
  \caption{Threading}\label{fig:main-threading}
\end{subfigure}\hfill
\begin{subfigure}[t]{0.195\textwidth}
  \includegraphics[width=\linewidth]{figs/placeholder_main_coffee.pdf}
  \caption{Coffee}\label{fig:main-coffee}
\end{subfigure}\hfill
\begin{subfigure}[t]{0.195\textwidth}
  \includegraphics[width=\linewidth]{figs/placeholder_main_square.pdf}
  \caption{Square}\label{fig:main-square}
\end{subfigure}\hfill
\begin{subfigure}[t]{0.195\textwidth}
  \includegraphics[width=\linewidth]{figs/placeholder_main_tpa.pdf}
  \caption{ThreePieceAssembly~[TPA]}\label{fig:main-tpa}
\end{subfigure}\hfill
\begin{subfigure}[t]{0.195\textwidth}
  \includegraphics[width=\linewidth]{figs/placeholder_main_locomotion.pdf}
  \caption{Locomotion (IsaacLab)}\label{fig:main-locomotion}
\end{subfigure}
% TODO(fig): put ONE shared legend above or below the row (e.g. a legend-only
% subfigure spanning \textwidth), never five per-panel legends.
\caption{%
  \textbf{Main sim-to-sim result: adaptation under a matched target budget.}
  Target performance vs.\ target transitions consumed, one panel per task; the
  $x$-axis ends at the budget cap of Def.~\ref{def:budget} ([TODO: $100$k]
  transitions). All arms start from the identical source assets and receive the
  identical $n \le 20$ target demonstrations; none sees target reward shaping,
  the target physics parameters, or a target expert. The grey dash-dotted line
  is the frozen source expert (zero-shot), i.e.\ the level every curve starts
  from, and the black dotted line is the same planner run with \emph{oracle}
  target dynamics --- the ceiling attainable once model error is removed, shown
  as a reference and not as a competitor. [TODO: the two zero-interaction
  baselines, target BC and a domain-randomized source expert, appear as
  additional horizontal lines.] Shading is $\pm 1$ seed spread over [TODO: $k$]
  seeds; each point is an independent re-evaluation over [TODO: $\ge 300$]
  episodes.
  \textbf{SimOPD is the only arm that both starts high and keeps improving
  within the budget}, reaching $\ph{0.--}$ on average and closing $\ph{--\%}$
  of the reference gap; the model-free arms are still near their zero-shot level
  when the budget runs out.
  % TODO(fig): the locomotion panel is a DIFFERENT regime -- dense reward,
  % periodic contact, no terminal goal predicate -- so Eq.~\eqref{eq:sparse}
  % does not hold there. Say so in the caption in one clause and use it as a
  % generality claim ("the ordering survives outside the sparse-reward
  % contact-rich regime we analyzed"), NOT as another instance of the analysis.
  % Its y-axis is normalized return, not success rate: label it separately and
  % do not let a shared y-label imply the units are the same.
  % TODO(fig): if the locomotion budget is not comparable in transitions
  % (locomotion envs are massively parallel in IsaacLab), plot it against
  % environment steps AND state the wall-clock/env-count in the caption, or the
  % "matched budget" claim silently breaks on that panel.
}
\label{fig:main-sim2sim}
\end{figure*}

% TODO(float 6b, optional): a single summary panel -- "fraction of the reference
% gap (pi_s -> oracle planner) closed at the budget", one grouped bar per task.
% This is NOT decoration: several tasks sit near the ceiling, where an additive
% delta understates the effect, and the curves alone will make the strongest
% tasks look like null results. Either append it as a 6th panel of
% Fig.~\ref{fig:main-sim2sim} or fold it into the ablation float below.

[TODO: text for Fig.~\ref{fig:main-sim2sim}. Read it in three passes, in this
order, so the figure does the argument rather than merely reporting it.
\emph{(i) Families, not curves.} The three line styles are the three ways of
consuming interaction data (Sec.~\ref{preliminary}), and inside the budget they
are ordered: every planning curve lies above every Dyna curve, which lies above
every model-free curve, on all [TODO: five] panels. The ordering, not the margin
of any single arm, is the result.
\emph{(ii) Shape, not endpoint.} Point out what only a curve can show: the
model-free arms are still flat at the cap --- they have not begun to convert
interaction into performance --- while the model-based arms have already spent
their budget productively. This is the same statement as
Sec.~\ref{sec:mf-vs-mb}, now made at the level of the final system.
\emph{(iii) Ours vs.\ the runner-up and vs.\ the ceiling.} SimOPD improves on
SimDist by $\ph{+0.--}$ averaged over tasks; attribute this to the two
components and forward-reference the ablation. On [TODO: which] tasks the SimOPD
curve enters the noise band of the oracle-dynamics reference before the cap ---
i.e.\ the remaining model error has stopped costing anything, the strongest form
the claim of Sec.~\ref{sec:what-to-adapt} can take. On [TODO: which] tasks a gap
remains; say honestly what it is (tolerance-window width, or residual persistent
model bias) rather than leaving it unremarked.]

% TODO(float 7, ablation -- to be added once the method text is written):
% component ablation of SimOPD, same panel layout as Fig.~\ref{fig:main-sim2sim}
% but arms = { SimDist base, + ensemble only, + optimistic exploration only,
% + pessimistic planning only, + anchored BC proposal only, full SimOPD }.
% Curves again, for the same reason: the anchored-BC component acts on the
% STARTING level and the ensemble on the SLOPE/ceiling, and an endpoint bar
% chart cannot distinguish those two mechanisms. Keep it a separate float from
% Fig.~\ref{fig:main-sim2sim}: the main figure answers "is it best", the
% ablation answers "why", and merging them makes both unreadable.
% Fewer tasks are acceptable here (2-3 panels) if space is tight -- but say
% which and why, and do not silently pick the ones where it looks best.
% Two things this ablation must control for, both of which have bitten us:
%   - planner hyper-parameters do not transfer across arms; each row needs its
%     own re-tuned init_std / base_policy_std or the ablation measures tuning.
%   - the anchor coefficient's benefit is target-dependent (>20x spread across
%     targets), so report per-task and never average it into one number.
